\documentclass[12pt,oneside,paper=5.5in:8.5in,DIV=11,chapterprefix=false,headings=big]{scrbook}

\usepackage{kidstyle}
\usepackage{graphicx}

% What a Logo program drew, from figures/, generated by tools/check-programs.sh.
% No extension: most figures are vector PDFs, a few are screenshots.
\newcommand{\drawing}[2][1.25in]{%
  \begin{center}
    \includegraphics[height=#1,width=0.7\linewidth,keepaspectratio]{figures/#2}
  \end{center}}

\begin{document}

\frontmatter

% ---------------------------------------------------------------------------
% Cover
% ---------------------------------------------------------------------------
% art/cover.jpg takes over the middle of the cover. Without it the cover falls
% back to drawing its own picture, so the book still builds on a bare checkout.

\begin{titlepage}
\centering

\vspace*{\stretch{0.45}}

\begin{tikzpicture}
  \foreach \i/\c in {0/kidred, 1/kidorange, 2/kidgreen, 3/kidblue, 4/kidpurple, 5/kidteal}
    \fill[\c] (\i*0.62,0) circle (0.17);
\end{tikzpicture}

\vspace{1.2em}

{\fontsize{34}{38}\selectfont\bfseries\color{kidink} The Kid\\[0.1em]
 \color{kidblue} Computer\par}

\vspace{0.8em}

{\Large\bfseries\color{kidink!70} Christmas 2025\par}

\vspace{1.0em}

\IfFileExists{art/cover.jpg}{%
  % Trimmed to the ink: the generated file has a wide white margin all round.
  % The file is 2048px tagged 300dpi, so LaTeX sees it as 491bp square and the
  % trim is in those bp, not pixels. Re-measure if the art is ever re-exported.
  \includegraphics[trim=35bp 111bp 69bp 105bp,clip,width=0.80\textwidth]%
                  {art/cover.jpg}%
}{%
  % Placeholder art: a friendly machine with the turtle on the screen.
  \begin{tikzpicture}[scale=1.0,line join=round]
    \definecolor{caseA}{HTML}{F0E6D2}
    \definecolor{caseB}{HTML}{CDBFA3}
    % monitor case
    \fill[caseA,rounded corners=10pt] (-2.5,-0.2) rectangle (2.5,3.2);
    \draw[caseB,line width=2pt,rounded corners=10pt] (-2.5,-0.2) rectangle (2.5,3.2);
    % screen
    \fill[kidink,rounded corners=6pt] (-2.05,0.35) rectangle (2.05,2.85);
    % the turtle, part way through drawing a star
    \definecolor{screenink}{HTML}{FFC24B}
    \definecolor{screenturtle}{HTML}{5BC46E}
    \begin{scope}[shift={(0,1.55)}]
      \draw[screenink,line width=1.8pt,line cap=round]
        (90:0.9) -- (234:0.9) -- (18:0.9) -- (162:0.9) -- (306:0.9) -- cycle;
      \begin{scope}[shift={(18:1.2)},scale=0.55]
        \kidturtlebody{screenturtle}{kidgreen}{kidgreendark}{white}
      \end{scope}
    \end{scope}
    % power light
    \fill[kidgreen] (2.15,0.08) circle (0.09);
    % stand
    \fill[caseB] (-0.55,-0.85) rectangle (0.55,-0.2);
    \fill[caseB,rounded corners=4pt] (-1.35,-1.15) rectangle (1.35,-0.8);
    % keyboard
    \fill[caseA,rounded corners=5pt] (-2.2,-2.15) rectangle (2.2,-1.35);
    \draw[caseB,line width=1.6pt,rounded corners=5pt] (-2.2,-2.15) rectangle (2.2,-1.35);
    \foreach \x in {-1.95,-1.65,...,1.75}
      \fill[caseB] (\x,-1.75) rectangle ++(0.2,0.22);
    \fill[caseB] (-0.75,-2.05) rectangle (0.95,-1.87);
  \end{tikzpicture}%
}

\vspace{1.0em}

{\Large\bfseries\color{kidink} Sally \& Penny Utterback\par}

\vspace*{\stretch{1}}

{\footnotesize\color{kidink!60} Updated \today\par}
\end{titlepage}

% ---------------------------------------------------------------------------

\renewcommand*{\contentsname}{What Is Inside}
\chaptericon{\faListUl}
\kidcompactheadtrue
\tableofcontents
\kidcompactheadfalse

\mainmatter

% ---------------------------------------------------------------------------
\chaptericon{\faDesktop}
\chapter{Welcome}

\lede{This is a computer for doing a few things carefully.}

You can make pictures.
You can move a turtle.
You can type small programs.

When you are done with something,
you can always come back.

% ---------------------------------------------------------------------------
\chaptericon{\faBars}
\chapter{The Menu}

\lede{When the computer turns on, you see the menu.}

The menu lets you choose what to do.

Use the arrow keys to pick what you want.
Press \key{Enter} to start it.

To leave and go back to the menu, hold down
\key{Ctrl} and \key{Alt}, then press \key{Backspace}.

% ---------------------------------------------------------------------------
\chaptericon{\faPaintBrush}
\chapter{Drawing}

\lede{Pick ``Draw Pictures'' from the menu.}

You get a big white page
and rows of buttons around the edge.

Pick a colour.
Pick a brush.
Then hold the mouse button down
and move the mouse across the page.

There are stamps you can press onto the page,
and lines and shapes you can stretch,
and an eraser for taking things away.

If you do something you did not mean to,
press Undo and it goes away again.

Press Save to keep a picture.
Press Open to look at the pictures you have kept.

% ---------------------------------------------------------------------------
\chaptericon{\faGlobeAmericas}
\chapter{Websites}

\lede{Your computer can talk to other computers.}

It sends a message out along a wire.
That wire goes to another wire,
and after a great many wires
the message arrives at a computer far away.
All of those computers and wires together
are called the Internet.

Some of those far-away computers keep pages
for other people to read.
A page can have words on it, or pictures, or a game.
When a lot of pages belong together,
we call them a website.

Every page has an address,
the way a house has an address.
Your computer uses the address to ask
the far-away computer for the page,
and that computer sends the page back.
It does this every time you look at a page,
which is why a website only works
when your computer can reach the Internet.

\section{The Websites on This Computer}

Some things on the menu are websites.

The address is already filled in for you,
and there is nowhere to type a different one,
so you cannot end up somewhere else by accident.

You can only visit the websites we picked together.

\begin{heads}[If the page says it is blocked]
Nothing is broken. It just means that page is not one of ours.
Press \key{Esc} to go back to the page you were reading.
\end{heads}

If you want a new website added,
ask, and we can talk about it.

% ---------------------------------------------------------------------------
\chaptericon{\faBookOpen}
\chapter{Looking Up Words}

\lede{Pick ``Look Up a Word'' from the menu.}

Type a word and press \key{Enter}.
The computer tells you what it means.

If you are not sure how to spell it,
type your best guess anyway.
The computer will show you words
that look close to what you typed,
and one of them is probably the one you wanted.

Press \key{Enter} on an empty line to go back to the menu.

The words come from people all over the world
who wrote them down in simple English
so that anyone learning to read can understand them.

% ---------------------------------------------------------------------------
\chaptericon{\faClock[regular]}
\chapter{The Clock}

\lede{Pick ``Clock'' from the menu.}

There are two clocks on the screen,
and they always say the same thing.

One is round, with the numbers 1 to 12 around the edge
and three hands in the middle.
The short fat hand points at the hour.
The long hand counts the minutes.
The thin red one goes all the way round
once every minute, so you can watch it move.

The other clock writes the time out in numbers,
and the day and the date underneath.

If you are still learning the hands,
read the round clock first
and then look at the numbers to see if you were right.

% ---------------------------------------------------------------------------
\chaptericon{\faHourglassHalf}
\chapter{The Timer}

\lede{Pick ``Timer'' from the menu.}

The computer asks how many minutes you want.
Type a number and press \key{Enter}.

Then the big numbers count backwards,
and music plays while they count.

When it gets to zero the music stops,
the screen says TIME'S UP,
and a bell rings until somebody stops it.

\begin{heads}[To stop the timer]
Press \key{Ctrl} and \key{C} at the same time.
This works while it is counting
and while the bell is ringing.
\end{heads}

% ---------------------------------------------------------------------------
\chaptericon{\faStopwatch}
\chapter{The Stopwatch}

\lede{Pick ``Stopwatch'' from the menu.}

The timer counts backwards.
The stopwatch counts forwards.

It starts at zero the moment it opens,
and it keeps going until you stop it.

Use it to find out how long something takes.
How long can you stand on one foot?
How long does it take you to put your shoes on?

To stop it, press \key{Ctrl} and \key{C}.

% ---------------------------------------------------------------------------
\chaptericon{\kidturtlew[0.72]}
\chapter[The Turtle (Logo)]{The Turtle (Logo): Telling the Computer to Draw}

\lede{Instead of drawing yourself, you can tell the computer what to draw,
so that it happens exactly like you want.}

The turtle listens to words.

You can tell the turtle to move.
You can tell it to turn.

Try typing:

\begin{code}
FD 50
RT 90
LT 90
\end{code}

\begin{heads}[If the turtle does not do it]
The turtle knows only a few words,
and it wants them spelled exactly right.

\chip{FD 50} works. \chip{FD50} does not,
because the word and the number
need a space between them.

When the turtle does not understand,
it prints a message instead of drawing.
Nothing is broken and you have not hurt anything.
Look at what you typed, fix the one letter
or the one space that is wrong,
and press \key{Enter} again.
\end{heads}

\section{Square}

\begin{code}
FD 100
RT 90
FD 100
RT 90
FD 100
RT 90
FD 100
RT 90
\end{code}

\drawing{square-long}

That is a lot of typing.
You can also say:

\begin{code}
REPEAT 4 [FD 100 RT 90]
\end{code}

\begin{trythis}
\begin{itemize}
\item \chip{100} to make it bigger or smaller
\item \chip{90} to see what happens
\end{itemize}
\end{trythis}

\section{Triangle}

\begin{code}
REPEAT 3 [FD 100 RT 120]
\end{code}

\drawing{triangle}

\section{Circle}

\begin{code}
REPEAT 36 [FD 10 RT 10]
\end{code}

\drawing{thirty-six-sides}

\begin{trythis}[Try this]
\begin{itemize}
\item changing the numbers
\item making it repeat more times
\end{itemize}
\end{trythis}

\section{Star}

\begin{code}
REPEAT 5 [FD 100 RT 144]
\end{code}

\drawing{star}

\section{Flower}

\begin{code}
REPEAT 8 [REPEAT 4 [FD 50 RT 90] RT 45]
\end{code}

\drawing[1.5in]{flower-eight}

\section{Useful Commands}

\begin{commandlist}
\cmd{CS}{clear the screen}
\cmd{HOME}{send the turtle back to the middle}
\cmd{PU}{pen up (move without drawing)}
\cmd{PD}{pen down (draw again)}
\end{commandlist}

% ---------------------------------------------------------------------------
% The contents needs two pages at this size; break it where the book does,
% so the turtle is one page and BASIC the other rather than splitting a chapter.
\addtocontents{toc}{\protect\newpage}
\chaptericon{\faTerminal}
\chapter[BASIC]{BASIC: Telling the Computer What to Do}

\lede{You can do Daddy's job! This computer can run small programs.}

Programs are made from lines with numbers.

You type the lines.
Then you type:

\begin{code}
RUN
\end{code}

\begin{heads}[If the computer does not understand a line]
Computers are very fussy.
A missing quote mark,
or a word spelled almost right,
is enough to stop one.

When that happens the computer prints a message
and tells you which line it did not like.
That is the computer trying to help,
not the computer being cross.

Type that line again with the same number at the front.
The new line takes the place of the old one.
Then type \chip{RUN} again.
\end{heads}

\section{Useful Commands}

\begin{commandlist}
\cmd{RUN}{makes the program go}
\cmd{LIST}{shows what you typed}
\cmd{NEW}{erase and start over}
\cmd{CLS}{clear the screen}
\end{commandlist}

\begin{heads}[If the screen fills up or will not stop]
Press \key{Ctrl} and \key{C} at the same time.
\end{heads}

\section{Growing Stars}

This makes stars grow down the screen.

\begin{code}
10 CLS
20 FOR I = 1 TO 20
30 PRINT STRING$(I, "*")
40 NEXT I
50 END
\end{code}

\begin{trythis}
\begin{itemize}
\item \chip{20} to a bigger number
\item \chip{"*"} to another letter
\end{itemize}
\end{trythis}

\section{Letter Rain}

Letters fall forever!

\begin{code}
10 CLS
20 PRINT CHR$(65 + RND * 26);
30 GOTO 20
\end{code}

To stop it, press \key{Ctrl} and \key{C}.

\section{Bouncing Word}

The word moves back and forth.

\begin{code}
10 CLS
20 FOR X = 1 TO 20
30 PRINT SPC(X) + "HI!"
40 NEXT X
50 FOR X = 20 TO 1 STEP -1
60 PRINT SPC(X) + "HI!"
70 NEXT X
80 GOTO 10
\end{code}

\begin{trythis}
\begin{itemize}
\item \chip{"HI!"} to your name
\item the numbers to make it move more
\end{itemize}
\end{trythis}

\section{Number Explosion}

Numbers fill the screen.

\begin{code}
10 CLS
20 FOR I = 1 TO 10
30 PRINT I; I; I; I; I
40 NEXT I
50 END
\end{code}

\begin{trythis}[Try this]
\begin{itemize}
\item printing more numbers
\item changing \chip{I} to \chip{I * I}
\end{itemize}
\end{trythis}

\section{Checkerboard}

This makes a pattern.

\begin{code}
10 CLS
20 FOR R = 1 TO 10
30 PRINT "* * * * * "
40 PRINT " * * * * *"
50 NEXT R
60 END
\end{code}

\begin{trythis}[Try this]
\begin{itemize}
\item changing the symbols
\item adding more rows
\end{itemize}
\end{trythis}

\section{Name Shout}

The computer asks your name!

\begin{code}
10 CLS
20 INPUT "WHAT IS YOUR NAME"; N$
30 FOR I = 1 TO 10
40 PRINT N$
50 NEXT I
60 END
\end{code}

\begin{trythis}[Try this]
\begin{itemize}
\item changing how many times it prints
\item printing something else
\end{itemize}
\end{trythis}

% ---------------------------------------------------------------------------
% ---------------------------------------------------------------------------
% One page holding every word the rest of the book teaches, so a kid who has
% forgotten one does not have to hunt through the chapter it came from. It has
% to stay one page to be worth having, which is why the headings here are
% lightweight rather than \section and the two word lists sit side by side.
\newcommand{\cardhead}[1]{%
  \par\vspace{0.4em}{\bfseries\large\color{kidaccent}#1}\par\vspace{0.2em}}
\newenvironment{cardlist}[1][1.6cm]
  {\commandlist[leftmargin=#1,labelwidth=\dimexpr#1-0.2cm\relax,
                labelsep=0.2cm,itemsep=0.05em,topsep=0.2em]}
  {\endcommandlist}

\chaptericon{\faKey}
\kidcompactheadtrue
\chapter{What to Type}
\kidcompactheadfalse

\cardhead{Getting Around}

\begin{cardlist}[3.6cm]
\cmd{Enter}{start the thing you picked}
\cmd{Esc}{go back, on a website}
\cmd{Ctrl + C}{stop something that will not stop}
\cmd{Ctrl + Alt + Backspace}{go back to the menu}
\end{cardlist}

\noindent
\begin{minipage}[t]{0.48\linewidth}
\RaggedRight
\cardhead{The Turtle}
\begin{cardlist}
\cmd{FD 50}{go forward}
\cmd{RT 90}{turn right}
\cmd{LT 90}{turn left}
\cmd{REPEAT}{do it again}
\cmd{CS}{clear the screen}
\cmd{HOME}{go to the middle}
\cmd{PU}{pen up}
\cmd{PD}{pen down}
\end{cardlist}
\end{minipage}\hfill
\begin{minipage}[t]{0.48\linewidth}
\RaggedRight
\cardhead{BASIC}
\begin{cardlist}
\cmd{RUN}{make it go}
\cmd{LIST}{show your lines}
\cmd{NEW}{start over}
\cmd{CLS}{clear the screen}
\cmd{PRINT}{show something}
\cmd{INPUT}{ask and wait}
\cmd{FOR}{do lines again}
\cmd{GOTO}{jump to a line}
\cmd{END}{stop here}
\end{cardlist}
\end{minipage}

% Pad to a multiple of 4 pages so the booklet folds cleanly. These are useful
% pages rather than blank ones, so padding costs nothing.

\dotgridpage{Pictures for the Turtle}
\dotgridpage{My Programs}
\dotgridpage{My Ideas}

\end{document}
